\documentclass[a4paper,11pt]{article}

\usepackage{amsmath,amssymb,amsfonts}
\usepackage{fontspec}
\usepackage{libertinus-otf}
\usepackage[french]{babel}
\usepackage{array}
\usepackage{booktabs}
\usepackage{longtable}
\usepackage{geometry}
\usepackage{multirow}
\usepackage[table]{xcolor}
\usepackage{colortbl}
\usepackage{fancyhdr}
\usepackage{lastpage}
\usepackage{enumitem}
\usepackage{eurosym}
\usepackage{tikz}
\usepackage[most]{tcolorbox}

% --- Géométrie calibrée pour exactement 4 pages ---
\geometry{
    a4paper,
    left=1.1cm,
    right=1.1cm,
    top=1.6cm,
    bottom=1.5cm,
    headheight=20pt,
    headsep=10pt,
    footskip=20pt
}

\setlength{\parindent}{0pt}
\setlength{\tabcolsep}{4.5pt}
\renewcommand{\arraystretch}{1.18}
\setlength{\aboverulesep}{1.5pt}
\setlength{\belowrulesep}{1.5pt}

% --- Palette institutionnelle harmonisée ---
\definecolor{bleuprimaire}{RGB}{20, 52, 98}
\definecolor{bleusecondaire}{RGB}{31, 95, 168}
\definecolor{bleufond}{RGB}{238, 244, 251}
\definecolor{grisaccent}{RGB}{100, 116, 139}
\definecolor{vertvrai}{RGB}{16, 120, 60}
\definecolor{rougefaux}{RGB}{195, 20, 20}
\definecolor{filetgris}{RGB}{220, 226, 235}

% --- Configuration TikZ pour des tableaux parfaits et compacts ---
\tikzset{
    tabcadre/.style={draw=bleuprimaire, line width=0.8pt},
    tabligne/.style={draw=bleuprimaire!45, line width=0.6pt},
    tabrep/.style={
        draw=rougefaux!90,
        fill=rougefaux!8,
        text=rougefaux,
        font=\bfseries\footnotesize,
        rounded corners=2.5pt,
        inner sep=1.8pt,
        line width=0.6pt
    },
    tabfleche/.style={
        ->,
        >=stealth,
        thick,
        draw=bleusecondaire,
        shorten <= 3pt,
        shorten >= 3pt
    }
}

% --- Pilules de points de barème ---
\newcommand{\pts}[1]{%
    \tcbox[on line, size=small, colback=bleufond, colframe=bleusecondaire!50,
           arc=2.5mm, boxrule=0.6pt, top=1.2pt, bottom=1.2pt, left=4pt, right=4pt]{%
        \small\bfseries\textcolor{bleuprimaire}{#1}%
    }%
}

% --- Cadre de réponse cible pour le texte hors-tableau ---
\newcommand{\repbox}[1]{%
    \tcbox[on line, size=small, colback=rougefaux!6, colframe=rougefaux!90,
           arc=1.2mm, boxrule=0.7pt, top=1pt, bottom=1pt, left=3.5pt, right=3.5pt]{%
        \bfseries\textcolor{rougefaux}{#1}%
    }%
}

\newcommand{\vrai}{\textcolor{vertvrai}{\textbf{[VRAI]}}}
\newcommand{\faux}{\textcolor{rougefaux}{\textbf{[FAUX]}}}
\newcommand{\subpt}[1]{\textcolor{bleusecondaire}{\bfseries [#1]}}
\newcommand{\tol}[1]{\textcolor{grisaccent}{\footnotesize\textit{(#1)}}}

% --- En-têtes et pieds de page ---
\pagestyle{fancy}
\fancyhf{}
\lhead{\small\bfseries\textcolor{bleuprimaire}{BAC BLANC} \textcolor{grisaccent}{-- Terminale Spécialité Mathématiques}}
\rhead{\small\bfseries\textcolor{bleusecondaire}{Barème de correction -- Sujet 1}}
\cfoot{\small\textcolor{grisaccent}{Page \thepage\ sur \pageref{LastPage}}}

\renewcommand{\headrulewidth}{0.5pt}
\renewcommand{\headrule}{\hbox to\headwidth{\color{bleusecondaire!60}\leaders\hrule height \headrulewidth\hfill}}
\renewcommand{\footrulewidth}{0.3pt}
\renewcommand{\footrule}{\hbox to\headwidth{\color{filetgris}\leaders\hrule height \footrulewidth\hfill}}

% --- SQUELETTE : marqueur magenta des zones à remplir (utilisable en texte ET en mode maths) ---
\newcommand{\ph}[1]{\ifmmode\text{\color{magenta}[#1]}\else\textcolor{magenta}{[#1]}\fi}


% =========================================================================
\begin{document}

% -------------------------------------------------------------------------
% EN-TÊTE DU BARÈME
% -------------------------------------------------------------------------
\begin{tcolorbox}[
    enhanced,
    colback=bleufond,
    colframe=bleuprimaire,
    arc=1.6mm,
    boxrule=0.8pt,
    top=4pt, bottom=4pt, left=8pt, right=8pt
]
    \textbf{\Large\textcolor{bleuprimaire}{Grille d'évaluation chiffrée}} \hfill \textbf{\Large\textcolor{bleuprimaire}{Total : 20 points}} \\[0.15em]
    \textbf{\textcolor{bleusecondaire}{Sujet 1 (Jour 1)}} \hspace{0.2cm}\textcolor{grisaccent}{\textbullet}\hspace{0.2cm} \textcolor{grisaccent}{Bac Blanc Mathématiques} \hfill \small\textcolor{grisaccent}{Durée : 4 heures \textbullet\ Calculatrice autorisée}
\end{tcolorbox}

\vspace{0.10cm}

% -------------------------------------------------------------------------
% TABLEAU DE NOTATION EN CONTINU
% -------------------------------------------------------------------------
\begin{longtable}{
    >{\centering\arraybackslash\bfseries}p{1.1cm}
    >{\raggedright\arraybackslash\bfseries}p{4.1cm}
    !{\color{filetgris}\vrule width 0.7pt}
    >{\raggedright\arraybackslash}p{11.1cm}
    >{\centering\arraybackslash}p{1.2cm}
}

\toprule[1.2pt]
\rowcolor{bleuprimaire}
\textcolor{white}{\bfseries Réf.} &
\textcolor{white}{\bfseries Compétence / Question} &
\textcolor{white}{\bfseries Critères d'évaluation \& visuels attendus} &
\textcolor{white}{\bfseries Barème} \\
\midrule[1pt]
\endfirsthead

\toprule[1.2pt]
\rowcolor{bleuprimaire}
\textcolor{white}{\bfseries Réf.} &
\textcolor{white}{\bfseries Compétence / Question} &
\textcolor{white}{\bfseries Critères d'évaluation \& visuels attendus} &
\textcolor{white}{\bfseries Barème} \\
\midrule[1pt]
\endhead

% ==================== EXERCICE 1 ====================
\rowcolor{bleusecondaire!15}
\multicolumn{4}{l}{
    \bfseries\textcolor{bleuprimaire}{\large Exercice 1 \hspace{0.2cm}\textbullet\hspace{0.2cm} \ph{Titre de l'exercice 1}} \hfill \bfseries 4 points
} \\* \midrule[0.6pt]

\textbf{Q.1} &
\ph{Compétence Q.1} &
\ph{VRAI ou FAUX}\newline
\textbullet\ \subpt{0,50} \ph{Critère 1}.\newline
\textbullet\ \subpt{0,25} \ph{Critère 2}.\newline
\textbullet\ \subpt{0,25} \ph{Critère 3} : \repbox{\ph{réponse attendue}}. &
\pts{1,00} \\
\cmidrule(lr){2-4}
\textbf{Q.2} &
\ph{Compétence Q.2} &
\ph{VRAI ou FAUX}\newline
\textbullet\ \subpt{0,25} \ph{Critère 1}.\newline
\textbullet\ \subpt{0,25} \ph{Critère 2}.\newline
\textbullet\ \subpt{0,50} \ph{Critère 3} : \repbox{\ph{réponse attendue}}.\newline
\tol{\ph{Tolérance / méthode alternative acceptée}}. &
\pts{1,00} \\
\cmidrule(lr){2-4}
\textbf{Q.3} &
\ph{Compétence Q.3} &
\ph{VRAI ou FAUX}\newline
\textbullet\ \subpt{0,50} \ph{Critère 1}.\newline
\textbullet\ \subpt{0,50} \ph{Critère 2} : \repbox{\ph{réponse attendue}}. &
\pts{1,00} \\
\cmidrule(lr){2-4}
\textbf{Q.4} &
\ph{Compétence Q.4} &
\ph{VRAI ou FAUX}\newline
\textbullet\ \subpt{0,50} \ph{Critère 1}.\newline
\textbullet\ \subpt{0,50} \ph{Critère 2} : \repbox{\ph{réponse attendue}}.\newline
\tol{\ph{Tolérance / méthode alternative acceptée}}. &
\pts{1,00} \\


\midrule[1pt]

% ==================== EXERCICE 2 ====================
\rowcolor{bleusecondaire!15}
\multicolumn{4}{l}{
    \bfseries\textcolor{bleuprimaire}{\large Exercice 2 \hspace{0.2cm}\textbullet\hspace{0.2cm} \ph{Titre de l'exercice 2}} \hfill \bfseries 5 points
} \\* \midrule[0.6pt]

% --- PARTIE A ---
\rowcolor{bleusecondaire!10}
\multicolumn{4}{l}{%
    \bfseries\textcolor{bleuprimaire}{\textsc{Partie A}} \;\textcolor{grisaccent}{\textbullet}\; \textcolor{bleuprimaire}{\ph{Titre de la partie A}} \hfill \textcolor{bleusecondaire}{\bfseries 2,50 points}%
} \\* \midrule[0.5pt]

\textbf{A.1} &
\ph{Compétence A.1} &
\vspace{-0.20cm}
\begin{center}
\begin{tikzpicture}[xscale=1.75, yscale=0.50, font=\small, baseline=(current bounding box.center)]
    \node (R) at (0,0) {};
    \node (A) at (1.5,1.2) {$\dots$};
    \node (B) at (1.5,-1.2) {$\dots$};
    \node (AD) at (3.2,1.8) {$\dots$};
    \node (ADb) at (3.2,0.6) {$\dots$};
    \node (BD) at (3.2,-0.6) {$\dots$};
    \node (BDb) at (3.2,-1.8) {$\dots$};

    \draw[thick, draw=bleuprimaire!80] (R) -- (A) node[midway, above left=-1.5pt] {\footnotesize$\dots$};
    \draw[thick, draw=bleuprimaire!80] (R) -- (B) node[midway, below left=-1.5pt] {\footnotesize$\dots$};
    \draw[thick, draw=bleuprimaire!80] (A) -- (AD) node[midway, above] {\footnotesize$\dots$};
    \draw[thick, draw=bleuprimaire!80] (A) -- (ADb) node[midway, below] {\footnotesize$\dots$};
    \draw[thick, draw=bleuprimaire!80] (B) -- (BD) node[midway, above] {\footnotesize$\dots$};
    \draw[thick, draw=bleuprimaire!80] (B) -- (BDb) node[midway, below] {\footnotesize$\dots$};
\end{tikzpicture}
\end{center}
\vspace{-0.20cm}
\ph{Description de la figure attendue}. &
\pts{0,50} \\
\cmidrule(lr){2-4}
\textbf{A.2} &
\ph{Compétence A.2} &
\textbullet\ \subpt{0,25} \ph{Critère 1}.\newline
\textbullet\ \subpt{0,50} \ph{Critère 2} : \repbox{\ph{réponse attendue}}. &
\pts{0,75} \\
\cmidrule(lr){2-4}
\textbf{A.3} &
\ph{Compétence A.3} &
\ph{Critère} : \repbox{\ph{réponse attendue}}. &
\pts{0,50} \\
\cmidrule(lr){2-4}
\textbf{A.4} &
\ph{Compétence A.4} &
\textbullet\ \subpt{0,25} \ph{Critère 1}.\newline
\textbullet\ \subpt{0,50} \ph{Critère 2} : \repbox{\ph{réponse attendue}}. &
\pts{0,75} \\


\midrule[0.6pt]

% --- PARTIE B ---
\rowcolor{bleusecondaire!10}
\multicolumn{4}{l}{%
    \bfseries\textcolor{bleuprimaire}{\textsc{Partie B}} \;\textcolor{grisaccent}{\textbullet}\; \textcolor{bleuprimaire}{\ph{Titre de la partie B}} \hfill \textcolor{bleusecondaire}{\bfseries 2,50 points}%
} \\* \midrule[0.5pt]

\textbf{B.1.a} &
\ph{Compétence B.1.a} &
\ph{Critère} : \repbox{\ph{réponse attendue}}. &
\pts{0,50} \\
\cmidrule(lr){2-4}
\textbf{B.1.b} &
\ph{Compétence B.1.b} &
\textbullet\ \subpt{0,25} \ph{Critère 1}.\newline
\textbullet\ \subpt{0,25} \ph{Critère 2} : \repbox{\ph{réponse attendue}}. &
\pts{0,50} \\
\cmidrule(lr){2-4}
\textbf{B.1.c} &
\ph{Compétence B.1.c} &
\ph{Critère} : \repbox{\ph{réponse attendue}}. &
\pts{0,50} \\
\cmidrule(lr){2-4}
\textbf{B.2.a} &
\ph{Compétence B.2.a} &
\ph{Critère} : \repbox{\ph{réponse attendue}}. &
\pts{0,50} \\
\cmidrule(lr){2-4}
\textbf{B.2.b} &
\ph{Compétence B.2.b} &
\textbullet\ \subpt{0,25} \ph{Critère 1}.\newline
\textbullet\ \subpt{0,25} \ph{Critère 2} : \repbox{\ph{réponse attendue}}.\newline
\tol{\ph{Tolérance / méthode alternative acceptée}}. &
\pts{0,50} \\


\midrule[1pt]

% ==================== EXERCICE 3 ====================
\rowcolor{bleusecondaire!15}
\multicolumn{4}{l}{
    \bfseries\textcolor{bleuprimaire}{\large Exercice 3 \hspace{0.2cm}\textbullet\hspace{0.2cm} \ph{Titre de l'exercice 3}} \hfill \bfseries 6 points
} \\* \midrule[0.6pt]

% --- PARTIE A ---
\rowcolor{bleusecondaire!10}
\multicolumn{4}{l}{%
    \bfseries\textcolor{bleuprimaire}{\textsc{Partie A}} \;\textcolor{grisaccent}{\textbullet}\; \textcolor{bleuprimaire}{\ph{Titre de la partie A}} \hfill \textcolor{bleusecondaire}{\bfseries 2,50 points}%
} \\* \midrule[0.5pt]

\textbf{A.1} &
\ph{Compétence A.1} &
\textbullet\ \subpt{0,25} \ph{Critère 1}.\newline
\textbullet\ \subpt{0,25} \ph{Critère 2} : \repbox{\ph{réponse attendue}}. &
\pts{0,50} \\
\cmidrule(lr){2-4}
\textbf{A.2} &
\ph{Compétence A.2} &
\textbullet\ \subpt{0,25} \ph{Critère 1}.\newline
\textbullet\ \subpt{0,50} \ph{Critère 2} :\newline
-- \ph{Étape 1}.\newline
-- \ph{Étape 2}.\newline
-- Conclusion : \repbox{\ph{réponse attendue}}. &
\pts{0,75} \\
\cmidrule(lr){2-4}
\textbf{A.3.a} &
\ph{Compétence A.3.a} &
\ph{Critère} : \repbox{\ph{réponse attendue}}. &
\pts{0,50} \\
\cmidrule(lr){2-4}
\textbf{A.3.b} &
\ph{Compétence A.3.b} &
\ph{Critère} : \repbox{\ph{réponse attendue}}. &
\pts{0,25} \\
\cmidrule(lr){2-4}
\textbf{A.4} &
\ph{Compétence A.4} &
\ph{Critère} : \repbox{\ph{réponse attendue}}. &
\pts{0,50} \\


\midrule[0.6pt]

% --- PARTIE B ---
\rowcolor{bleusecondaire!10}
\multicolumn{4}{l}{%
    \bfseries\textcolor{bleuprimaire}{\textsc{Partie B}} \;\textcolor{grisaccent}{\textbullet}\; \textcolor{bleuprimaire}{\ph{Titre de la partie B}} \hfill \textcolor{bleusecondaire}{\bfseries 3,50 points}%
} \\* \midrule[0.5pt]

\textbf{B.1} &
\ph{Compétence B.1} &
\ph{Critère 1}.\newline
\ph{Critère 2} : \repbox{\ph{réponse attendue}}. &
\pts{0,50} \\
\cmidrule(lr){2-4}
\textbf{B.2} &
\ph{Compétence B.2} &
\ph{Critère} : \repbox{\ph{réponse attendue}}. &
\pts{0,25} \\
\cmidrule(lr){2-4}
\textbf{B.3.a} &
\ph{Compétence B.3.a} &
\ph{Critère} : \repbox{\ph{réponse attendue}}. &
\pts{0,25} \\
\cmidrule(lr){2-4}
\textbf{B.3.b} &
\ph{Compétence B.3.b} &
\ph{Critère} : \repbox{\ph{réponse attendue}}. &
\pts{0,50} \\
\cmidrule(lr){2-4}
\textbf{B.3.c} &
\ph{Compétence B.3.c} &
\ph{Critère} : \repbox{\ph{réponse attendue}}. &
\pts{0,50} \\
\cmidrule(lr){2-4}
\textbf{B.4} &
\ph{Compétence B.4} &
\textbullet\ \subpt{0,25} \ph{Critère 1}.\newline
\textbullet\ \subpt{0,25} \ph{Critère 2} : \repbox{\ph{réponse attendue}}. &
\pts{0,50} \\
\cmidrule(lr){2-4}
\textbf{B.5.a} &
\ph{Compétence B.5.a} &
\textbullet\ \subpt{0,25} \ph{Critère 1}.\newline
\textbullet\ \subpt{0,25} \ph{Critère 2} : \repbox{\ph{réponse attendue}}.\newline
\tol{\ph{Tolérance / méthode alternative acceptée}}. &
\pts{0,50} \\
\cmidrule(lr){2-4}
\textbf{B.5.b} &
\ph{Compétence B.5.b} &
\textbullet\ \subpt{0,25} \ph{Critère 1}.\newline
\textbullet\ \subpt{0,25} \ph{Critère 2} : \repbox{\ph{réponse attendue}}. &
\pts{0,50} \\


\midrule[1pt]

% ==================== EXERCICE 4 ====================
\rowcolor{bleusecondaire!15}
\multicolumn{4}{l}{
    \bfseries\textcolor{bleuprimaire}{\large Exercice 4 \hspace{0.2cm}\textbullet\hspace{0.2cm} \ph{Titre de l'exercice 4}} \hfill \bfseries 5 points
} \\* \midrule[0.6pt]

% --- PARTIE A ---
\rowcolor{bleusecondaire!10}
\multicolumn{4}{l}{%
    \bfseries\textcolor{bleuprimaire}{\textsc{Partie A}} \;\textcolor{grisaccent}{\textbullet}\; \textcolor{bleuprimaire}{\ph{Titre de la partie A}} \hfill \textcolor{bleusecondaire}{\bfseries 2,00 points}%
} \\* \midrule[0.5pt]

\textbf{A.1} &
\ph{Compétence A.1} &
\textbullet\ \subpt{0,25} \ph{Critère 1}.\newline
\textbullet\ \subpt{0,25} \ph{Critère 2} : \repbox{\ph{réponse attendue}}. &
\pts{0,50} \\
\cmidrule(lr){2-4}
\textbf{A.2} &
\ph{Compétence A.2} &
\vspace{-0.18cm}
\begin{center}
\begin{tikzpicture}[xscale=1.0, yscale=0.56, font=\footnotesize, baseline=(current bounding box.center)]
    \draw[tabcadre] (0,0) rectangle (5.8, -3.1);
    \draw[tabligne] (0,-0.75) -- (5.8,-0.75);
    \draw[tabligne] (0,-1.45) -- (5.8,-1.45);
    \draw[tabcadre] (1.3,0) -- (1.3,-3.1);

    \node at (0.65, -0.37) {\bfseries $x$};
    \node at (0.65, -1.10) {\bfseries $g'(x)$};
    \node at (0.65, -2.25) {\bfseries $g(x)$};

    \node at (1.8, -0.37) {$\dots$};
    \node at (5.2, -0.37) {$\dots$};

    \node at (3.5, -1.10) {$\dots$};

    \node[tabrep] (v1) at (1.8, -2.75) {$\dots$};
    \node (v2) at (5.2, -1.85) {$\dots$};

    \draw[tabfleche] (v1.north east) -- (v2.south west);
\end{tikzpicture}
\end{center}
\vspace{-0.18cm}
\ph{Critère} : \repbox{\ph{réponse attendue}}. &
\pts{0,50} \\
\cmidrule(lr){2-4}
\textbf{A.3} &
\ph{Compétence A.3} &
\ph{Critère} : \repbox{\ph{réponse attendue}}. &
\pts{0,50} \\
\cmidrule(lr){2-4}
\textbf{A.4} &
\ph{Compétence A.4} &
\vspace{-0.18cm}
\begin{center}
\begin{tikzpicture}[xscale=1.0, yscale=0.52, font=\footnotesize, baseline=(current bounding box.center)]
    \draw[tabcadre] (0,0) rectangle (5.8, -1.6);
    \draw[tabligne] (0,-0.8) -- (5.8,-0.8);
    \draw[tabcadre] (1.4,0) -- (1.4,-1.6);

    \node at (0.7, -0.4) {\bfseries $x$};
    \node at (0.7, -1.2) {\bfseries $g(x)$};

    \node at (2.2, -0.4) {$\dots$};
    \node at (3.8, -0.4) {$\dots$};
    \node at (5.3, -0.4) {$\dots$};

    \node[tabrep] at (3.0, -1.2) {$\dots$};
    \node at (3.8, -1.2) {$\dots$};
    \node[tabrep] at (4.6, -1.2) {$\dots$};
\end{tikzpicture}
\end{center}
\vspace{-0.18cm}
\ph{Critère}. &
\pts{0,50} \\


\midrule[0.6pt]

% --- PARTIE B ---
\rowcolor{bleusecondaire!10}
\multicolumn{4}{l}{%
    \bfseries\textcolor{bleuprimaire}{\textsc{Partie B}} \;\textcolor{grisaccent}{\textbullet}\; \textcolor{bleuprimaire}{\ph{Titre de la partie B}} \hfill \textcolor{bleusecondaire}{\bfseries 3,00 points}%
} \\* \midrule[0.5pt]

\textbf{B.1} &
\ph{Compétence B.1} &
\textbullet\ \subpt{0,50} \ph{Critère 1}.\newline
\textbullet\ \subpt{0,25} \ph{Critère 2} : \repbox{\ph{réponse attendue}}. &
\pts{0,75} \\
\cmidrule(lr){2-4}
\textbf{B.2} &
\ph{Compétence B.2} &
\textbullet\ \subpt{0,25} \ph{Critère 1}.\newline
\textbullet\ \subpt{0,25} \ph{Critère 2} : \repbox{\ph{réponse attendue}}.\newline
\tol{\ph{Tolérance / méthode alternative acceptée}}. &
\pts{0,50} \\
\cmidrule(lr){2-4}
\textbf{B.3} &
\ph{Compétence B.3} &
\vspace{-0.18cm}
\begin{center}
\begin{tikzpicture}[xscale=1.0, yscale=0.56, font=\footnotesize, baseline=(current bounding box.center)]
    \draw[tabcadre] (0,0) rectangle (6.8, -3.1);
    \draw[tabligne] (0,-0.75) -- (6.8,-0.75);
    \draw[tabligne] (0,-1.45) -- (6.8,-1.45);
    \draw[tabcadre] (1.3,0) -- (1.3,-3.1);

    \node at (0.65, -0.37) {\bfseries $x$};
    \node at (0.65, -1.10) {\bfseries $f'(x)$};
    \node at (0.65, -2.25) {\bfseries $f(x)$};

    \node at (1.8, -0.37) {$\dots$};
    \node at (4.0, -0.37) {$\dots$};
    \node at (6.2, -0.37) {$\dots$};

    \node[tabrep] at (2.9, -1.10) {$\dots$};
    \node at (4.0, -1.10) {$\dots$};
    \node[tabrep] at (5.1, -1.10) {$\dots$};

    \node[tabrep] (botL) at (1.8, -2.75) {$\dots$};
    \node[tabrep] (topM) at (4.0, -1.85) {$\dots$};
    \node[tabrep] (botR) at (6.2, -2.75) {$\dots$};

    \draw[tabfleche] (botL.north east) -- (topM.south west);
    \draw[tabfleche] (topM.south east) -- (botR.north west);
\end{tikzpicture}
\end{center}
\vspace{-0.18cm}
\textbullet\ \subpt{0,25} \ph{Critère 1}.\newline
\textbullet\ \subpt{0,25} \ph{Critère 2}. &
\pts{0,50} \\
\cmidrule(lr){2-4}
\textbf{B.4} &
\ph{Compétence B.4} &
\textbullet\ \subpt{0,25} \ph{Critère 1}.\newline
\textbullet\ \subpt{0,25} \ph{Critère 2} : \repbox{\ph{réponse attendue}}. &
\pts{0,50} \\
\cmidrule(lr){2-4}
\textbf{B.5.a} &
\ph{Compétence B.5.a} &
\ph{Critère} : \repbox{\ph{réponse attendue}}. &
\pts{0,25} \\
\cmidrule(lr){2-4}
\textbf{B.5.b} &
\ph{Compétence B.5.b} &
\ph{Critère} : \repbox{\ph{réponse attendue}}. &
\pts{0,50} \\*


\midrule[1.2pt]
\rowcolor{bleuprimaire}
\multicolumn{3}{r}{\textcolor{white}{\bfseries\large TOTAL GÉNÉRAL}} &
\textcolor{white}{\bfseries\large 20 pts} \\
\bottomrule[1.2pt]

\end{longtable}

\vspace{0.08cm}

% -------------------------------------------------------------------------
% CONSIGNES COMMUNES DE CORRECTION (PAGE 4)
% -------------------------------------------------------------------------
\begin{tcolorbox}[
    enhanced,
    colback=white,
    colframe=bleuprimaire,
    boxrule=0.7pt,
    leftrule=3.5pt,
    arc=1.5mm,
    title={\bfseries Consignes communes d'harmonisation (Règles d'or)},
    coltitle=bleuprimaire,
    colbacktitle=bleufond,
    titlerule=0.3pt,
    top=4pt, bottom=4pt, left=8pt, right=8pt
]
\small
\begin{itemize}[leftmargin=*, label=\textcolor{bleusecondaire}{\small\textbullet}, itemsep=8pt, topsep=1pt, parsep=0pt]
    \item \textbf{\ph{Intitulé de la règle 1}} : \ph{Texte de la règle 1}.
    \item \textbf{\ph{Intitulé de la règle 2}} : \ph{Texte de la règle 2}.
    \item \textbf{\ph{Intitulé de la règle 3}} : \ph{Texte de la règle 3}.
    \item \textbf{\ph{Intitulé de la règle 4}} : \ph{Texte de la règle 4}.
    \item \textbf{\ph{Intitulé de la règle 5}} : \ph{Texte de la règle 5}.
    \item \textbf{\ph{Intitulé de la règle 6}} : \ph{Texte de la règle 6}.
    \item \textbf{\ph{Intitulé de la règle 7}} : \ph{Texte de la règle 7}.
\end{itemize}
\end{tcolorbox}

\end{document}